\bookinterlude{III}{IV}{Quem fica com a chave}
Até sexta-feira o telefone tocava depois que alguém já tinha escolhido o torno. Eu consultava a floresta, contava quantas árvores apostavam numa peça dentro da medida e passava o palpite adiante. Na sexta o telefone toca mais cedo. É o encarregado com a barra na mão. “Em qual dos dois eu ponho?”

No torno da janela, nove em cada dez peças têm saído na medida. O da porta voltou do conserto ontem. O cliente vem buscar uma peça hoje e encomendou mais cem iguais para segunda. Vamos usar a mesma regulagem nos dois dias.

Abro uma página no caderno e desenho duas colunas, uma para cada torno. Se eu mandar a barra para o da janela, de tarde saberei se ele acertou mais uma. A coluna do torno da porta vai continuar em branco. Posso chamar mais dez pessoas para dar palpite. Para escrever uma medida naquela coluna, alguém vai ter de gastar uma barra lá.

Deixo a chave com você. Onde vai pôr a barra de sexta? Enquanto você pensa, o telefone toca de novo. O cliente cancelou as cem peças de segunda. Quanto você ainda pagaria para descobrir hoje como ficou o torno que voltou do conserto?
